\section{Introduction}
\label{sec:introduction}

Device and browser fingerprints occupy an uneasy position in Web security. They can provide a passive signal for associating a request with a previously observed environment, augmenting passwords or informing risk-based authentication (RBA) decisions~\cite{alaca2016device,papaioannou2023survey}. At the same time, mobile fingerprints may be insufficiently distinctive or stable to serve as a trusted-device factor on their own~\cite{spooren2015harmful}, and the same attributes can be used for unwanted tracking~\cite{laperdrix2020browser,w3c2025fingerprinting}. A security system therefore cannot treat a long feature vector, or even a highly distinctive value, as an integrity proof. It must reason about where each observation came from, whether it was available under the current platform conditions, which relationships are expected to hold, and which deviations admit benign explanations.

The adversarial setting makes this requirement more acute. Browser-facing attributes can be altered through injected scripts, browser configuration, debugging interfaces, automation frameworks, or purpose-built evasion services. Gummy Browsers demonstrates that a targeted adversary can capture a victim fingerprint and orchestrate another browser to imitate it~\cite{liu2022gummy}. Large-scale measurements of commercial traffic likewise show that adversarial and benign fingerprints form a meaningful problem for deployed Web defenses~\cite{wu2023manyfaces}. Yet manipulating isolated attributes is not equivalent to producing a coherent environment. FP-Scanner showed that anti-fingerprinting countermeasures often introduce contradictory combinations of browser attributes~\cite{vastel2018fpscanner}; more recently, FP-Inconsistent found spatial and temporal inconsistencies in evasive bot traffic and learned rules that expose part of this manipulation~\cite{venugopalan2025fpinconsistent}. BrowserFM further demonstrates that browser attributes can be linked systematically to the configurations and environments that generate them~\cite{huyghe2025browserfm}. These studies establish consistency as an important signal, but they primarily reason \emph{within} a browser fingerprint or between a browser and its generating configuration.

Android hybrid applications expose a broader boundary. The host application can observe Native build, hardware, display, memory, sensor, installation, and debugging state. The WebView host contributes provider, package, JavaScript-bridge, user-agent, and security-configuration information. JavaScript inside the embedded page sees a browser-like runtime containing navigator, screen, viewport, Canvas, WebGL, timezone, resource, and automation-facing properties. These surfaces are technically distinct but operationally coupled. Android applications can customize WebView behavior and create explicit Java--JavaScript channels; recent large-scale studies show that these channels are widely used to connect native identifiers, Web APIs, tracking code, and network destinations~\cite{weerasekara2025tracking,datta2026crossboundary}. Earlier systems such as BridgeTaint and iWanDroid demonstrate that security-relevant evidence must often be tracked across the Java--JavaScript boundary rather than analyzed on one side alone~\cite{bai2019bridgetaint,tiwari2023demand}.

This paper asks a different question from information-flow or tracking measurement: \emph{can the multiple observations produced during one controlled Android hybrid-app session be treated as provenance-bound integrity evidence, and which explicit relationships change when part of the environment is manipulated?} A Native Android 13 observation paired with a desktop Windows user agent is not merely an unusual browser string; it contradicts the host surface that produced the Web runtime. A WebView provider version and a runtime Chromium version may be comparable only at a coarse major-version level. A physical-resolution relation requires orientation, insets, viewport, and density tolerance. A missing Browser capture is not a vector of zeroes, and an unparseable user agent is not evidence that two versions match. These distinctions motivate HybridGuard.

HybridGuard is a research framework for fixed-contract collection, provenance-aware pairing, deterministic evidence extraction, source-separated relation execution, and auditable reporting. Its current in-app collection unit, \appfp, contains 84 Android Native signals, 26 WebView-host signals, and 67 in-app Web-runtime signals. Each field is accompanied by a state such as observed, unsupported, permission-denied, runtime-error, timeout, or not-applicable. Accepted payloads are bound to collection receipts, batches, release identifiers, payload hashes, and quality-control records. Deterministic extractors convert raw fields into an evidence bundle containing parsed facts, cross-layer comparisons, applicability conditions, and explicit unassessed states. Relations derived from official Android/Chrome field semantics are maintained separately from empirically mined predicates; official documents ground what a field means, but the resulting risk relation remains a project inference.

The current controlled experiment uses \appfp only. This is a deliberate evidence boundary rather than an assertion that the external browser is unimportant: the repository contains an independent \browserfp probe and a provenance-complete \pairedfp path, but complete external-browser observations are currently too sparse for a defensible attack-effect evaluation. We therefore evaluate verified two-state \emph{baseline $\rightarrow$ attack-active} pairs over the three in-app layers, report exactly which fields changed, which relations became violated, and which configurations remain uncovered. Once sufficiently complete paired data are collected, the same protocol will be rerun to compare \appfp-only, \browserfp-only where meaningful, and joint \pairedfp variants.

On the current frozen cohort, 229 normal observations and 69 qualified controlled pairs are available. Nine executable semantic relations grounded in official field definitions yield 51 baseline-no-alert to active-alert transitions; ten independently maintained device-mined predicates yield 33. Thirty-three pairs trigger both sources, eighteen trigger only the official-derived relation set, and eighteen trigger neither. The uncovered pairs include manipulations restricted to WebDriver, resource values, languages, plugins/MIME types, timezone, and screen metrics. These values characterize the behavior of frozen rules on frozen controlled configurations. They are \emph{not} an estimate of population-level recall, benign false-positive rate, unknown-attack generalization, or production RBA effectiveness.

Taken together, this draft makes five connected contributions. First, it introduces a provenance-aware multi-surface data model that combines fixed App and external-browser contracts with per-field state semantics, receipt/batch/hash binding, provenance-complete pairing, and honest App-only retention when browser capture is incomplete. The collection and pairing path is already implemented, although large-scale paired attack data remain future work. Second, it converts \appfp payloads into deterministic Native--WebView--Web evidence whose relations explicitly encode applicability, tolerance, counterexamples, and not-assessed outcomes rather than treating missing data as normal or directly concatenating all raw fields. Third, it separates project-derived semantic relations grounded in official field definitions from device-mined predicates, executes the two sources independently, and retains the field-level evidence that supports each decision. Fourth, it evaluates these frozen relation sets on verified baseline-to-active pairs supplied by the controlled-manipulation workstream and reports both observed alert transitions and concrete coverage gaps without relabeling the counts as attack recall. Finally, it establishes an implemented \browserfp and \pairedfp data path while reserving any claim about the incremental detection value of the external-browser surface for a later complete-data ablation.

The remainder of this draft first positions HybridGuard relative to device-fingerprint authentication, browser inconsistency detection, and Android WebView analysis (Section~\ref{sec:background-related}). It then defines the observation surfaces and threat model (Section~\ref{sec:problem}), presents the system and evidence pipeline (Sections~\ref{sec:system} and~\ref{sec:evidence}), specifies the two-state experimental methodology (Section~\ref{sec:methodology}), reports the current \appfp results (Section~\ref{sec:current-results}), and separates historical preliminary evidence from the future \browserfp study (Sections~\ref{sec:historical} and~\ref{sec:browser-future}).
